\cotefile{33}
\scanpage{18}
à des $f$ de type fini au lieu de '$f$' propre. Se ramène au cas d'une \uline{imm. ouverte}. Se résoud par l'\unsure{affirmative|affirm.} lorsqu'on dispose de la résolution des singularités\inserted{$+$ pureté}, p. ex. au cas où \inserted{pour ``bons préschémas''}, \unsure{ou plus|ou plan|ou plutôt}\inserted{\unsure{bons présch|bien précisé}} \unsure{cf. le|cgt le|gén.} cas lorsqu'on disposera de résolution des singularités \deleted{\unsure{des|de}} dans ce cadre....

14) \uline{Les} $R^{i}_{!}f\unsure{{}_{*}|{}}(F)$.

Soit $f:X\longrightarrow Y$ morphisme compactifiable, i.e. qui se factorise en

\diagram{\begin{tikzcd}[ampersand replacement=\&,column sep=large]
X \arrow[r,dash,"i","{\mathrm{imm.\ ouv.}}"'] \& \bar X \arrow[r,"{\bar f}","{\mathrm{propre}}"'] \& Y
\end{tikzcd}}

(p. ex. $f$ qu. proj. et $\deleted{f}Y$ quasi-cpt\unsure{?|.}).

Pour tt faisceau de torsion $F$ sur $X$, considérons $i_!(F)$ [exact en $F$] puis les $R^{q}_{\deleted{\unsure{q|i}}}\bar f_*(i_!(F))$, forment un foncteur coh. en $F$. Il résulte de façon quasi-triviale du th. de changement de base que ce foncteur coh.

% archive sheet number: 17
